\section{Open Questions}
The following five questions are genuinely open given what this
replication (and the paper itself) exercises. Each is stated
with the basis (why it is not settled by the current evidence)
and concrete next steps that are feasible on free endpoints.

\subsection*{Q1. Strongly-correlated systems}
\textbf{Question.} Does the Q$^2$Chemistry platform maintain
accuracy on strongly-correlated systems (e.g.\ Cr$_2$, N$_2$ near
dissociation, transition-metal dimers) where UCCSD-family
ansatz\"e are known to break down?

\textbf{Basis.} The paper's H$_2$/Fig.~6 sanity check exercises a
trivially-complete regime (UCCSD spans the full CI space for
H$_2$/STO-3G, so machine-precision agreement is analytically
expected). The Cr$_2$ demo (C3) targets a strongly-correlated
system but only reports MPS wall-clock scaling on 768~cores, not
energy accuracy vs.\ an exact (e.g.\ DMRG) reference. Whether the
platform's ansatz\"e deliver chemical accuracy in the
multi-reference regime is not established.

\textbf{Next steps.} Run Q$^2$Chemistry on N$_2$ across the
dissociation curve (0.9--2.5~\AA) with UCCSD, $k$-UpCCGSD, and
ADAPT-VQE, comparing to a CASSCF(10,8)/DMRG reference. Free-endpoint
feasible via PySCF + OpenFermion (STO-3G/6-31G) for the
classical-simulator side and block2 (open-source DMRG) as the
reference. Report max energy error and non-parallelism error
(NPE) across the curve for each ansatz.

\subsection*{Q2. Integration with modern QEM}
\textbf{Question.} How does Q$^2$Chemistry's MPS backend interact
with modern quantum error mitigation (ZNE, PEC, virtual
distillation)? Does the platform expose hooks for QEM at the
circuit level, and how much of the reported near-FCI accuracy
survives a realistic noise model?

\textbf{Basis.} The paper presents noiseless MPS simulations
only; no reported integration with QEM techniques and no
benchmark under a hardware-realistic noise channel. Modern
quantum-chemistry-on-QC work (2024--2026 IBM/Quantinuum/IonQ) is
dominated by mitigation choices.

\textbf{Next steps.} Instrument an H$_2$/LiH VQE run through the
platform (or an OpenFermion + Mitiq stand-in if the Q$^2$Chemistry
source is not obtainable) with a depolarising + amplitude-damping
noise model calibrated to a modern superconducting device. Sweep
ZNE (Richardson, orders 1--3) and virtual distillation, and
report residual energy error vs.\ FCI as a function of the
two-qubit gate error rate. Baseline: Mitiq.

\subsection*{Q3. Catalyst / drug-molecule benchmarks}
\textbf{Question.} How does the platform perform on
catalytically-relevant or drug-relevant benchmark molecules
(FeMoco active site, cytochrome P450 iron centre, or a
small drug fragment such as caffeine)?

\textbf{Basis.} The paper's numerical demos are H$_2$ (Fig.~6)
and bulk silicon bands (Fig.~7). Neither is a target where
near-term quantum chemistry is pitched as delivering value.

\textbf{Next steps.} Select a canonical near-term-QC benchmark
set (Fe(II)-porphyrin CAS(32,34); Mn$_3$O$_4$ model cluster;
active-space Hamiltonians from the OpenFermion benchmark suite).
Run VQE-UCCSD and ADAPT-VQE at CAS-projected reduced qubit counts
(10--16 qubits, state-vector tractable). Compare energies to
DMRG/CASPT2 reference. Report parameter counts, iteration counts,
and energy error.

\subsection*{Q4. End-to-end noise robustness of the integrated
pipeline}
\textbf{Question.} How noise-robust is the integrated
Q$^2$Chemistry pipeline end-to-end (Hamiltonian generation
$\to$ VQE energy)? Do errors from finite-shot sampling, gate
noise, and readout error compound catastrophically or does the
platform's design absorb them gracefully?

\textbf{Basis.} Noiseless simulator results only; no analysis of
how the classical--quantum interface (Hamiltonian construction,
measurement grouping, parameter-update loop) behaves when the
quantum backend returns finite-shot noisy estimates.

\textbf{Next steps.} For H$_2$/STO-3G and LiH/STO-3G, run the VQE
loop with a shot-noise backend at
$N_{\text{shots}} \in \{10^2, 10^3, 10^4, 10^5\}$ and depolarising
gate errors in $\{0, 10^{-4}, 10^{-3}, 10^{-2}\}$. Report VQE
energy variance across 20 seeds per configuration, plus optimiser
convergence rate. Compare against a canonical Qiskit-Nature VQE
under the same noise model.

\subsection*{Q5. ML-driven ansatz selection}
\textbf{Question.} Would ML-driven ansatz selection (an RL agent
or a supervised model over molecular-graph features) reduce the
parameter count and circuit depth of Q$^2$Chemistry ansatz\"e
compared to hand-crafted symmetry-reduced UCCSD, without
sacrificing accuracy?

\textbf{Basis.} The paper reports 53 variational parameters for
its symmetry-reduced UCCSD on H$_2$/ccj-pVDZ --- substantial for a
two-electron system. ADAPT-VQE and other adaptive schemes prune
parameters heuristically; whether a learned model over molecule
descriptors can do better, and how it interacts with the
platform's ansatz-builder interface, is unexplored.

\textbf{Next steps.} Train a small RL agent (PPO with a
graph-neural-network state encoder over molecule + operator pool)
to select excitation operators one at a time, with accuracy-vs-FCI
as reward and circuit depth as penalty. Baselines: hand-crafted
UCCSD, standard ADAPT-VQE. Molecule set: H$_2$, LiH, BeH$_2$,
H$_4$-chain (STO-3G, tractable). Report the Pareto front of
(parameter count, energy error) and compare to the paper's
53-parameter figure.
